\documentclass[11pt]{article}
\usepackage[T1]{fontenc}
\usepackage{lmodern}
\usepackage[margin=1in]{geometry}
\usepackage{microtype}
\usepackage{amsmath,amssymb,amsthm,mathtools}
\usepackage{booktabs,longtable,array,tabularx}
\usepackage{enumitem}
\usepackage{xcolor}
\usepackage{listings}
\usepackage{fancyhdr}
\usepackage{hyperref}
\usepackage[nameinlink,capitalize,noabbrev]{cleveref}
\usepackage{bookmark}

\hypersetup{
  colorlinks=true,
  linkcolor=blue!55!black,
  citecolor=blue!55!black,
  urlcolor=blue!55!black,
  pdftitle={Median Mathematics in Intuitionistic Primitive Recursive Arithmetic},
  pdfauthor={OpenAI, prepared for Vladimir Reshetnikov}
}
\setlength{\parindent}{0pt}
\setlength{\parskip}{0.55em}
\setlist{nosep,leftmargin=2em}
\pagestyle{fancy}
\fancyhf{}
\fancyhead[L]{Median mathematics in iPRA}
\fancyhead[R]{16 November 2026}
\fancyfoot[C]{\thepage}
\setlength{\headheight}{14pt}

\lstset{
  basicstyle=\ttfamily\small,
  columns=fullflexible,
  breaklines=true,
  frame=single,
  backgroundcolor=\color{black!3},
  xleftmargin=0.5em,
  xrightmargin=0.5em
}

\newtheorem{theorem}{Theorem}[section]
\newtheorem{lemma}[theorem]{Lemma}
\newtheorem{corollary}[theorem]{Corollary}
\newtheorem{proposition}[theorem]{Proposition}
\newtheorem{definition}[theorem]{Definition}
\newtheorem{remark}[theorem]{Remark}
\newtheorem{example}[theorem]{Example}
\newtheorem{construction}[theorem]{Construction}

\newcommand{\iPRA}{\ensuremath{\mathrm{iPRA}}}
\newcommand{\Nat}{\mathbb N}
\newcommand{\Int}{\mathbb Z}
\newcommand{\Rat}{\mathbb Q}
\newcommand{\Real}{\mathbb R}
\newcommand{\Med}{\operatorname{Med}}
\newcommand{\card}{\operatorname{card}}
\newcommand{\abs}[1]{\left|#1\right|}
\newcommand{\norm}[1]{\left\lVert#1\right\rVert}
\newcommand{\floor}[1]{\left\lfloor#1\right\rfloor}
\newcommand{\ceil}[1]{\left\lceil#1\right\rceil}
\newcommand{\ind}{\mathbf 1}
\newcommand{\statusA}{\textbf{Finitary theorem}}
\newcommand{\statusB}{\textbf{Effective analytic schema}}
\newcommand{\statusC}{\textbf{Formalization route}}
\newcommand{\statusX}{\textbf{Correction}}

\title{\textbf{Median Mathematics in Intuitionistic\\Primitive Recursive Arithmetic}\\[0.5em]
\large Finitary proofs, effective analytic extensions, and a Lean formalization plan}
\author{Prepared for Vladimir Reshetnikov}
\date{16 November 2026}

\begin{document}
\maketitle

\begin{abstract}
This report reconstructs as much as possible of the mathematical content of the
English Wikipedia article \emph{Median} inside intuitionistic primitive
recursive arithmetic (\iPRA).  The finite core is unusually favorable to this
weak constructive setting.  A median of a nonempty finite sequence is defined
without division by two bounded counts.  For a sorted sequence
$x_0\leq\cdots\leq x_{n-1}$, the complete set of such medians is exactly
\[
  \left[x_{\floor{(n-1)/2}},\,x_{\floor{n/2}}\right].
\]
This theorem yields the odd and even formulas, existence, uniqueness,
quantile terminology, invariance under order transformations, robustness, and
the exact minimizers of total absolute deviation.  For even samples we work
with twice the conventional median, so no division is needed in the arithmetic
theory.

For finite weighted rational laws we prove the cumulative-distribution
characterization, existence, symmetry results, conditional-median optimality,
and exact odd-sample median distributions.  We give a denominator- and
square-root-free proof of the mean--median inequality
\[
  (2S-MW)^2\leq 4(WQ-S^2),
\]
where $W,S,Q$ are the total weight and first two raw moment numerators and $M$
is twice a median.  We also prove finite forms of Jensen's inequality for
medians, rank robustness, the replacement breakdown point, properties of
median filters, one-dimensional $k$-medians, Hodges--Lehmann and Theil--Sen
estimators, marginal medians, and the centerpoint theorem route.

Continuous distributions, logarithms, square roots, geometric medians, and
asymptotic normality are treated only through primitive-recursive regular
names and explicit moduli.  Several statements on the source page require
additional hypotheses; the report supplies corrections and exact finite
counterexamples.  The accompanying Lean files are explicit blueprints, not a
completed kernel-checked \texttt{Prf PRAAx} development.
\end{abstract}

\begin{center}
\fbox{\begin{minipage}{0.93\textwidth}
\textbf{Trust boundary.}  The arguments below are fully developed human
proofs.  ``Internal to \iPRA'' means that every object is finitely coded, every
algorithm is primitive recursive, and the indicated correctness proof can be
organized as quantifier-free induction.  It does \emph{not} mean that this
archive already contains completed Lean proof terms.  The Lean files contain
no \texttt{sorry} and introduce no theorem axiom, but they remain blueprints.
\end{minipage}}
\end{center}

\tableofcontents
\newpage

\section{Source scope and classification}

The target article ranges from elementary order statistics through robust
statistics, exact and asymptotic sampling distributions, multivariate
optimization, regression, filtering, clustering, and historical remarks
\cite{wiki-median}.  These topics have very different constructive strength.
We use four categories.

\begin{description}[style=nextline]
\item[Category A: finite exact mathematics.]  Lists, finite weighted laws,
bounded counts, sorting, finite sums, factorials, binomial coefficients, and
rational inequalities are primitive recursive.  Their proofs can be reduced
to quantifier-free induction.
\item[Category B: effective analysis.]  A real number or function is supplied
by rational interval approximations, together with primitive-recursive
precision, continuity, tail, and integrability moduli.  Claims become
$\varepsilon$--$N$ statements with computable witnesses.
\item[Category C: imported finite geometry or asymptotic theory.]  The
mathematical route is constructive, but a substantial library theorem---for
example fixed-dimensional Helly, effective resampling theory, or a strong
unimodal inequality---still has to be formalized.
\item[Category X: correction.]  Some informal sentences are false without
extra hypotheses.  We state and prove the corrected theorem, often with a
small counterexample to the unrestricted reading.
\end{description}

The separate \texttt{THEOREM\_LEDGER.md} records seventy article claims and
their status.

\section{The constructive arithmetic setting}

\subsection{What \iPRA supplies}

The ProveIt development presents primitive recursive arithmetic in
first-order intuitionistic logic.  Its language has a symbol for every
primitive-recursive function, defining equations for those symbols, equality,
successor axioms, and induction restricted to quantifier-free formulas
\cite{proveit-ipra}.  The checked derivation library already contains semiring
algebra, truncated subtraction, order, parity, bounded sums, pairing, search,
division, permutations, binomial coefficients, and related tools
\cite{proveit-derivations}.

A recurring formalization pattern is this.  A bounded loop is packaged as a
primitive-recursive function $h$.  The desired invariant is then a
quantifier-free equation or inequality involving $h$.  Induction on the loop
counter proves it.  Existential objects are not selected by an unbounded
choice principle: a sorting trace, index list, or bounded-search result is
computed explicitly and checked.

\subsection{Finite sequence coding}

A sequence code stores a length and a nested pairing of entries.  The
functions
\[
  \operatorname{len}(v),\qquad \operatorname{nth}(v,i),\qquad
  \operatorname{valid}(v)
\]
are total primitive-recursive functions; malformed inputs receive a harmless
default.  Bounded folds define counts and sums, for example
\[
 \operatorname{ltCount}(v,t)
   =\sum_{i<\operatorname{len}(v)}
      \ind[\operatorname{nth}(v,i)<t].
\]
The boolean comparison is decidable because arithmetic order is decidable.
The equation for the successor length is a direct bounded-sum recursion
identity.

The existing PA list-coding project already uses the right median design: a
nondecreasing permutation plus separate odd and even length equations for
\texttt{TwiceMedian}, avoiding division and parity-sensitive quotients
\cite{proveit-listcoding}.  We reuse that design, but the requested endpoint
is a new \iPRA\ derivation, not merely the existing PA semantic theorem.

\subsection{Integers and rationals}

Signed integers may be represented by pairs $(a,b)$ denoting $a-b$.  A
rational is a signed numerator and a positive denominator.  Equality and order
are tested by cross multiplication.  For a finite list of rationals one can
also multiply by the product $D$ of all positive denominators.  The map
$x\mapsto Dx$ is strictly increasing, so every order-statistic statement
reduces to signed integers.  At the end, division by $D$ is restored in the
external interpretation.

This scaling observation is important.  The finite median theory itself needs
no field axioms, no completed real line, and no choice.  Even the conventional
median is represented internally by its doubled value.

\section{The central finite theorem}

Let $x=(x_0,\ldots,x_{n-1})$ be a nonempty finite sequence in a linearly
ordered set with decidable comparison.  Let
\[
  y_0\leq y_1\leq\cdots\leq y_{n-1}
\]
be the output of an explicit insertion sort.  Define
\[
 \ell=\floor{\frac{n-1}{2}},\qquad
 u=\floor{\frac n2},\qquad L=y_\ell,\qquad U=y_u.
\]
For a candidate $t$, put
\[
 B_x(t)=\#\{i<n:x_i<t\},\qquad
 A_x(t)=\#\{i<n:t<x_i\}.
\]

\begin{definition}[Balance median]
A value $t$ is a balance median of $x$ if
\begin{equation}\label{eq:balance}
  2B_x(t)\leq n\quad\text{and}\quad 2A_x(t)\leq n.
\end{equation}
\end{definition}

This formulation is division-free and quantifier-free after the two counts
are made primitive-recursive symbols.

\begin{theorem}[Middle-interval theorem]\label{thm:middle-interval}
For every $t$,
\[
 t\text{ satisfies \eqref{eq:balance}}
 \quad\Longleftrightarrow\quad L\leq t\leq U.
\]
Consequently, the complete set of medians is the closed interval $[L,U]$.
\end{theorem}

\begin{proof}
Suppose first that $L\leq t\leq U$.  Because $t\leq y_u$, every entry strictly
below $t$ occurs before index $u$.  Hence $B_x(t)\leq u$, and
$2u\leq n$.  Similarly, because $y_\ell\leq t$, every entry strictly above
$t$ occurs after index $\ell$.  Thus
\[
 A_x(t)\leq n-\ell-1.
\]
A direct parity check gives $2(n-\ell-1)\leq n$.  Both balance inequalities
follow.

Conversely, assume $t<L$.  Every entry from index $\ell$ onward is strictly
above $t$, so
\[
 A_x(t)\geq n-\ell.
\]
For $n=2q+1$, this lower bound is $q+1$; for $n=2q$, it is also $q+1$.
In either case $2A_x(t)>n$, contradicting \eqref{eq:balance}.  The case $t>U$
is symmetric: the entries through index $u$ are strictly below $t$, so
$B_x(t)\geq u+1$ and $2B_x(t)>n$.
\end{proof}

Every inequality used in this proof is a polynomial or truncated-subtraction
identity in $n,q,\ell,u$.  The sortedness arguments are bounded.  Therefore
the proof has a direct quantifier-free-induction realization.

\begin{corollary}[Odd and even formulas]\label{cor:odd-even}
If $n=2q+1$, then $\ell=u=q$ and the median is uniquely $y_q$.  If $n=2q$,
then $\ell=q-1$ and $u=q$, and every $t\in[y_{q-1},y_q]$ is a median.  The
usual conventional choice is
\[
 m=\frac{y_{q-1}+y_q}{2}.
\]
Internally we define the unique doubled conventional median
\[
 M=2m=y_{q-1}+y_q.
\]
\end{corollary}

\begin{corollary}[Existence and uniqueness criterion]
Every nonempty finite sequence has a median.  Its median is unique exactly
when $L=U$; in particular every odd-length sequence has a unique median.
\end{corollary}

\begin{proof}
Insertion sort is primitive recursive and returns a sorted permutation.  The
value $L$ is therefore an explicit witness.  The uniqueness statement follows
from \cref{thm:middle-interval}.
\end{proof}

\subsection{Why sorting is constructively harmless}

Insertion sort is defined by primitive recursion on the input code.  Inserting
$a$ into a sorted list scans from the left until the first entry not below
$a$.  The correctness invariant has two finite components:

\begin{enumerate}
\item the result is nondecreasing;
\item a list of source indices witnesses that the result is a permutation of
the input.
\end{enumerate}

Both components are preserved by one insertion.  No appeal to ``there exists a
sorting permutation'' is needed.  The permutation certificate is useful later
because counts, sums, and products can be transported by a checked reindexing
lemma.

\section{Order, quantiles, and transformations}

\begin{theorem}[Order-only character]
The median predicate depends only on the comparison relation.  It is
independent of any distance or arithmetic structure.
\end{theorem}

\begin{proof}
The balance definition uses only $<$, finite counting, multiplication by two,
and comparison with the length.  The middle-interval characterization uses
only sorted order.  Therefore any finite ordinal scale---for example
$\mathrm F<\mathrm D<\mathrm C<\mathrm B<\mathrm A$---can be coded by ranks
and treated identically.
\end{proof}

For even ordinal data, the balance median is an interval between two adjacent
or separated grades.  A numerical midpoint is an extra convention, not an
order-theoretic object.

\begin{corollary}[Quantile names]
The median interval is the set of one-half quantiles.  The lower endpoint is
the lower $1/2$-quantile and the upper endpoint is the upper $1/2$-quantile.
Thus the same object is the second quartile, fifth decile, and fiftieth
percentile, once the lower/upper convention is fixed.
\end{corollary}

\begin{theorem}[Monotone transport]\label{thm:monotone}
Let $f$ be strictly increasing.  Then
\[
 t\in\Med(x)\quad\Longleftrightarrow\quad
 f(t)\in\Med(f(x_0),\ldots,f(x_{n-1})).
\]
If $f$ is strictly decreasing, the same conclusion holds, with the two count
inequalities exchanged.
\end{theorem}

\begin{proof}
For increasing $f$,
$x_i<t$ iff $f(x_i)<f(t)$ and $t<x_i$ iff $f(t)<f(x_i)$, so both counts are
unchanged.  For decreasing $f$, the first relation becomes
$f(t)<f(x_i)$ and the second becomes $f(x_i)<f(t)$; the balance conjunction is
symmetric in the two counts.
\end{proof}

\begin{remark}
The conventional midpoint of two even-sample middle values commutes with
\emph{affine} maps, including decreasing affine maps.  It does not commute
with an arbitrary nonlinear monotone map, although the full median interval
transports correctly.
\end{remark}

\section{Robustness and breakdown}

The phrase ``the median ignores extreme values'' admits an exact rank theorem.
Let $x$ and $z$ be samples of the same length, and suppose at most $r$
coordinates have been replaced.  Write $x_{(k)}$ and $z_{(k)}$ for their
zero-based order statistics.

\begin{theorem}[Rank perturbation]\label{thm:rank-perturbation}
For every $k<n$,
\begin{align*}
 r\leq k &\implies x_{(k-r)}\leq z_{(k)},\\
 k+r<n &\implies z_{(k)}\leq x_{(k+r)}.
\end{align*}
\end{theorem}

\begin{proof}
Assume $r\leq k$ and, toward a contradiction,
$z_{(k)}<x_{(k-r)}$.  At least $k+1$ entries of $z$ are then strictly below
$x_{(k-r)}$.  At most $r$ of them can be replacements, so at least
$k+1-r$ unchanged entries of $x$ are strictly below $x_{(k-r)}$.  But an order
statistic at index $k-r$ has at most $k-r$ entries strictly below it, a
contradiction.  The upper inequality follows by reversing the order.
\end{proof}

The proof is purely a bounded pigeonhole argument.  It does not assume the
values are numeric.

\begin{corollary}[Median rank envelope]
If fewer than half the observations are replaced, each middle order statistic
of the contaminated sample is trapped between two explicitly indexed order
statistics of the original sample.  In particular, changing values already
known to remain outside those ranks cannot change the median.
\end{corollary}

This explains censored observations.  If a small number of failures is known
only to exceed every completed time, their exact magnitudes are irrelevant as
long as they do not occupy a middle rank.

\begin{theorem}[Finite-sample replacement breakdown]\label{thm:breakdown}
For the conventional sample median of $n$ observations, the minimum number of
arbitrary replacements that can drive the estimate without bound is
$\ceil{n/2}$.  Hence the replacement breakdown fraction is
\[
  \frac{\ceil{n/2}}{n}.
\]
\end{theorem}

\begin{proof}
If $n=2q$, replace the largest $q$ chosen coordinates by a common value $H$.
The upper middle order statistic is then $H$, so the conventional median grows
like $H/2$.  With only $q-1$ replacements, \cref{thm:rank-perturbation} bounds
both middle order statistics by fixed original order statistics, so the
median cannot diverge.

If $n=2q+1$, replacing $q+1$ coordinates by $H$ makes the unique middle order
statistic equal to $H$.  With at most $q$ replacements, the rank perturbation
bound traps that order statistic between the original minimum and maximum.
\end{proof}

\section{Absolute-deviation optimality}

\subsection{Unweighted samples}

Define
\[
  D_x(t)=\sum_{i<n}\abs{x_i-t}.
\]

\begin{theorem}[Complete $L^1$ minimizer theorem]\label{thm:l1}
For a nonempty finite sample,
\[
 D_x(t)\leq D_x(a)\quad\text{for every }a
 \quad\Longleftrightarrow\quad t\in[L,U].
\]
Thus the minimizers of total, mean, or weighted-by-$1/n$ absolute deviation are
exactly the medians.
\end{theorem}

\begin{proof}
Sort the sample.  For $a\leq b$, the triangle inequality gives
\[
  \abs{a-t}+\abs{b-t}\geq b-a,
\]
with equality exactly when $a\leq t\leq b$.

If $n=2q$, pair $y_i$ with $y_{n-1-i}$ for $i<q$.  Summing the pairwise lower
bounds gives a constant lower bound for $D_x(t)$.  Equality in every pair
holds precisely when
\[
  t\in\bigcap_{i<q}[y_i,y_{n-1-i}]
   =[y_{q-1},y_q]=[L,U].
\]
If $n=2q+1$, the same $q$ outer pairs leave the central term
$\abs{y_q-t}$.  The outer-pair equalities permit an interval containing
$y_q$, and the central term is minimized uniquely at $t=y_q$.  Hence the
minimizer set is again $[L,U]$, now a singleton.
\end{proof}

The argument consists solely of finite pairing and an absolute-value case
split.  In the object theory, absolute value on signed integers is primitive
recursive, and the pair sum is a bounded fold.

\subsection{Weighted laws}

Let distinct support points $a_0<\cdots<a_{s-1}$ carry natural weights
$w_j$, with total $W=\sum_jw_j>0$.  Define
\[
 B(t)=\sum_{a_j<t}w_j,\qquad A(t)=\sum_{t<a_j}w_j,
 \qquad D(t)=\sum_jw_j\abs{a_j-t}.
\]
A weighted median satisfies $2B(t)\leq W$ and $2A(t)\leq W$.

\begin{theorem}[Weighted $L^1$ theorem]\label{thm:weighted-l1}
The weighted medians are exactly the minimizers of $D(t)$.
\end{theorem}

\begin{proof}
Expand the finite list by repeating $a_j$ exactly $w_j$ times.  This expansion
is primitive recursive and has length $W$.  Its ordinary counts and absolute
deviation sum are precisely $B,A,D$.  Apply \cref{thm:middle-interval,thm:l1}.
\end{proof}

An implementation need not materialize the expanded list.  It may instead use
weighted prefix sums.  Expansion is only a clean correctness proof.

\section{The mean--median inequality without square roots}

Let
\[
 W=\sum_jw_j,\qquad S=\sum_jw_ja_j,\qquad
 Q=\sum_jw_ja_j^2,
\]
and let $m$ be a weighted median.  Externally the mean and variance are
\[
 \mu=\frac SW,\qquad \sigma^2=\frac QW-\mu^2.
\]

\begin{theorem}[Cleared mean--median bound]\label{thm:mean-median}
For every weighted median $m$,
\begin{equation}\label{eq:mean-median-cleared}
  (S-Wm)^2\leq WQ-S^2.
\end{equation}
Equivalently, if $M=2m$ is a doubled median,
\begin{equation}\label{eq:mean-median-doubled}
  (2S-MW)^2\leq 4(WQ-S^2).
\end{equation}
Whenever division and square roots are interpreted, this is
$\abs{\mu-m}\leq\sigma$.
\end{theorem}

\begin{proof}
First,
\[
 \abs{S-Wm}
 =\abs{\sum_jw_j(a_j-m)}
 \leq\sum_jw_j\abs{a_j-m}.
\]
By \cref{thm:weighted-l1}, the right side is no larger than
$\sum_jw_j\abs{a_j-\mu}$.  Weighted Cauchy--Schwarz gives
\[
 \left(\sum_jw_j\abs{a_j-\mu}\right)^2
 \leq W\sum_jw_j(a_j-\mu)^2.
\]
No classical square-root theorem is needed for this step.  After clearing
denominators, Cauchy--Schwarz follows from the exact sum-of-squares identity
\[
 W\sum_jw_jd_j^2-\left(\sum_jw_j\abs{d_j}\right)^2
 =\sum_{i<j}w_iw_j\bigl(\abs{d_i}-\abs{d_j}\bigr)^2\geq0.
\]
Finally,
\[
 W\sum_jw_j(a_j-\mu)^2=WQ-S^2.
\]
Combining the three inequalities proves \eqref{eq:mean-median-cleared};
substituting $M=2m$ gives \eqref{eq:mean-median-doubled}.
\end{proof}

Every term in \eqref{eq:mean-median-doubled} is an integer or rational
numerator.  This is the preferred \iPRA\ statement.

\section{Finite probability distributions}

A finite rational law is a support list $a_0<\cdots<a_{s-1}$ with natural
weights $w_j$ and denominator $W$.  Its probability of an event is the event's
weight divided by $W$.  Since denominators are positive, all probability
inequalities can be cross-multiplied.

\begin{theorem}[CDF characterization]\label{thm:cdf}
Let
\[
 F(t)=\frac1W\sum_{a_j\leq t}w_j,
 \qquad F(t^-)=\frac1W\sum_{a_j<t}w_j.
\]
Then $t$ is a median iff
\[
  F(t^-)\leq\frac12\leq F(t).
\]
Equivalently,
$\Pr(X<t)\leq1/2$ and $\Pr(X>t)\leq1/2$.
\end{theorem}

\begin{proof}
The left inequality is $2B(t)\leq W$.  The right inequality is
$W\leq2(W-A(t))$, which simplifies to $2A(t)\leq W$.
\end{proof}

\begin{theorem}[Existence and median interval for a finite law]
Let $r$ be the least index with
$2\sum_{j\leq r}w_j\geq W$, and let $s$ be the greatest index with
$2\sum_{j\geq s}w_j\geq W$.  Then the complete median set is
$[a_r,a_s]$.
\end{theorem}

\begin{proof}
Minimality of $r$ gives strictly less than half the weight below $a_r$, while
the defining inequality gives at most half above it.  Thus $a_r$ is a median.
The symmetric argument proves that $a_s$ is a median.  Between the endpoints,
neither strict-side mass can increase past half.  Outside, one of the endpoint
arguments fails exactly as in \cref{thm:middle-interval}.
\end{proof}

The two endpoints differ exactly when cumulative mass is one half throughout
a gap.  Thus the familiar nonuniqueness of a distribution median is already a
finite rational phenomenon.

\section{Symmetry and named distributions}

\subsection{Finite symmetry}

A finite weighted law is symmetric about $c$ if its support admits an
involution $j\mapsto j^*$ such that
\[
 a_j+a_{j^*}=2c,\qquad w_j=w_{j^*}.
\]

\begin{theorem}[Symmetry center]\label{thm:symmetry}
For a finite symmetric law, $c$ is a median.  Its mean also equals $c$.  If
the law is given on a consecutive symmetric grid and its weights are unimodal,
the central grid point (or the two central grid points) belongs to the mode
plateau.  If a unique mode exists, it is the symmetry center and therefore the
grid must contain $c$.
\end{theorem}

\begin{proof}
The involution pairs every support point below $c$ with an equally weighted
point above $c$.  Hence the two strict-side masses are equal and each is at
most $W/2$, proving the median claim.  Pairwise cancellation gives
\[
 \sum_jw_j(a_j-c)=0,
\]
so the mean is $c$.  On a consecutive symmetric grid, symmetry carries every modal grid point to
another modal grid point.  Unimodality forces the modal set to be a central
block.  A unique mode makes that block a singleton fixed by reflection, hence
located at $c$.
\end{proof}

Finite uniform arithmetic progressions are immediate examples: their doubled
median and doubled mean are the sum of the endpoints.

\subsection{Effective real distributions}

For a continuous law we do not quantify over an arbitrary set-theoretic real
function.  We assume a regular CDF name: at precision $p$, primitive-recursive
functions return rational lower and upper bounds of width at most $2^{-p}$,
with monotonicity and a tail modulus.  A median name is a nested rational
interval whose CDF values are certified to straddle $1/2$.

Under that presentation, the standard formulas have constructive proofs:

\begin{itemize}
\item A symmetric normal or Cauchy law has its location parameter as median by
reflection of the CDF.  The Cauchy conclusion needs no mean.
\item For the uniform law on $[a,b]$, solving
$(m-a)/(b-a)=1/2$ gives $m=(a+b)/2$.
\item For an exponential law,
$1-e^{-\lambda m}=1/2$, hence $m=(\log2)/\lambda$.
\item For a Weibull law,
$1-e^{-(m/\lambda)^k}=1/2$, hence
$m=\lambda(\log2)^{1/k}$.
\item For the normalized power-law tail
$\Pr(X\geq x)=(x_{\min}/x)^{a-1}$, the half-tail equation gives
$m=2^{1/(a-1)}x_{\min}$.
\end{itemize}

In each transcendental case, ``equals'' means equality of regular names.  A
completed proof constructs logarithm, exponential, power, and root interval
evaluators and proves their defining inequalities.  The existing constructive
CLT packet supplies the appropriate regular-name architecture
\cite{proveit-clt}.

\section{Corrections concerning skew and unimodality}

\subsection{A finite skew counterexample}

The informal heuristic ``positive skew implies mean greater than median'' is
not a theorem.  The following counterexample is entirely finite.

\begin{proposition}\label{prop:skew-counterexample}
For the equally weighted sample
\[
  0,0,1,1,2,
\]
the mean is $4/5$, the median is $1$, and the third central moment is
$18/125>0$.
\end{proposition}

\begin{proof}
The mean is $(0+0+1+1+2)/5=4/5$, while the central order statistic is $1$.
Writing deviations from $4/5$ and cubing,
\[
 \frac15\left(2\left(-\frac45\right)^3
 +2\left(\frac15\right)^3
 +\left(\frac65\right)^3\right)
 =\frac{18}{125}>0.
\]
Thus even positive third central moment does not force mean above median.
\end{proof}

\subsection{The sharper unimodal bound needs hypotheses}

The article states a $\sqrt{3/5}$ improvement for unimodal distributions.  An
unrestricted discrete reading is false.  Put weight $2$ at $0$ and weight $3$
at $1$.  This law is unimodal, has median $1$, mean $3/5$, and variance
$6/25$.  Therefore
\[
 \frac{(m-\mu)^2}{\sigma^2}
 =\frac{4/25}{6/25}=\frac23>\frac35.
\]
The stronger bound must be tied to the precise continuous-unimodality
hypotheses of its source theorem.  Formalizing that version requires an
effective Gauss-type inequality and is classified as Category C.  The general
one-standard-deviation bound of \cref{thm:mean-median} is valid for all finite
laws.

Likewise, a statement about a monotonically decreasing \emph{density} is not
a statement about every finite decreasing probability table.  Its effective
continuous proof requires a regular density and an integrable-tail modulus.

\section{Jensen's inequality for medians}

A function $f$ on a linearly ordered domain is a C-function if every sublevel
set
\[
  \{x:f(x)\leq t\}
\]
is an interval (possibly empty or a singleton).  Convex functions are
C-functions, but convexity is stronger than necessary.

For a finite sample, let $U=x_{(\floor{n/2})}$ be the upper median, and let
$V$ be the upper median of the transformed sample $f(x_i)$.

\begin{theorem}[Finite Jensen inequality for medians]\label{thm:jensen-median}
If every relevant sublevel set of $f$ is an interval, then
\[
  f(U)\leq V.
\]
Equivalently, taking the suprema when median intervals are non-singletons,
\[
 f(\sup\Med(x))\leq \sup\Med(f(x)).
\]
\end{theorem}

\begin{proof}
Fix $t<f(U)$ and consider the interval
$C_t=\{x:f(x)\leq t\}$.  Since $U\notin C_t$, an interval $C_t$ lies wholly
below $U$ or wholly above $U$.  In the first case, it contains at most the
entries strictly below $U$; in the second, at most the entries strictly above
$U$.  Each count is at most $\floor{n/2}$.  Therefore no more than
$\floor{n/2}$ transformed observations are at most $t$.

But if the upper median $V$ were at most $t$, at least
$\floor{n/2}+1$ transformed observations would be at most $t$.  This is a
contradiction.  Hence $t<V$ for every $t<f(U)$, which in the finite rational
setting implies $f(U)\leq V$.
\end{proof}

For \iPRA, one does not quantify over arbitrary functions.  A finite table of
values $f(x_i)$ is accompanied by a certificate that, after sorting the
$x_i$, the indices satisfying $f(x_i)\leq t$ form a contiguous block for each
threshold drawn from the finite image.  The proof is then bounded counting.
Polynomial and piecewise-rational convex functions can obtain this certificate
from their usual monotonicity proof.

\section{Algorithms and complexity}

\subsection{Sorting versus selection}

Insertion sort gives the simplest correctness proof but costs quadratic time.
A median can be selected in linear comparison time by the deterministic
groups-of-five algorithm.

\begin{construction}[Median of medians]
Partition the input into groups of at most five, sort each group by a fixed
comparison network, recursively select the median of the group medians, use it
as a pivot, and recurse only into the side containing the desired rank.
\end{construction}

At least half of the group medians are no smaller than the pivot.  Except for
at most one incomplete group and the pivot group, each such median certifies
three elements no smaller than the pivot.  Symmetrically, at least roughly
$3n/10$ elements lie on each side.  With explicit ceiling terms, the cost
satisfies
\[
 T(n)\leq T(\ceil{n/5})+T(\ceil{7n/10}+6)+cn.
\]
Choose a constant $K$ large enough for the finitely many small cases and for
$K(1/5+7/10)n+K\cdot 8+cn\leq Kn$.  Strong induction then proves
$T(n)\leq Kn$.  The recursion tree, partition trace, and comparison count are
primitive-recursive; the inequality is a quantifier-free induction after the
cost function is defined.

A linear lower bound follows from an adversary argument: until an observation
has participated in a comparison chain, its possible rank cannot be excluded
from the middle.  Thus selection is $\Theta(n)$ in the comparison model.

\subsection{Sorting lower bound}

For distinct inputs, a correct comparison sort has at least $n!$ possible
output permutations.  A binary decision tree of depth $d$ has at most $2^d$
leaves, so $2^d\geq n!$.  This implication is finite and arithmetical.

For an explicit asymptotic lower bound, let $n=2m$.  The final $m$ factors of
$n!$ are at least $m$, hence $n!\geq m^m$.  If $m\geq2^r$, then
$n!\geq2^{mr}$ and therefore $d\geq mr$.  Taking
$r=\floor{\log_2m}$ gives the usual $\Omega(n\log n)$ bound without using real
logarithms inside the theory.

\subsection{Median of three and Tukey's ninther}

Median-of-three is a fixed compare-swap network and is therefore a
primitive-recursive exact median for three inputs.

\begin{theorem}[Nine-input ninther rank]\label{thm:ninther}
Divide nine distinct inputs into three triples, take each triple's median, and
then take the median of those three medians.  The output has global rank $4$,
$5$, or $6$ in one-based indexing.
\end{theorem}

\begin{proof}
The selected value is the median of its own triple, so that triple contributes
at least two values no greater than it.  One of the other triple medians is no
greater than it, and that triple contributes at least two more values no
greater than it.  Thus at least four inputs are no greater than the output.
The symmetric argument gives at least four no smaller.  Its rank is therefore
between $4$ and $6$.
\end{proof}

The exact checker exhausts all $9!$ permutations.  It finds ranks $4,5,6$ in
$77{,}760$, $207{,}360$, and $77{,}760$ cases respectively.

\subsection{Remedian resource bounds}

Fix an odd buffer size $b$.  Level zero stores incoming values until $b$ are
available, replaces them by their median, and promotes that value to level
one; levels repeat the rule.  Processing $n$ inputs performs at most
\[
 n+n/b+n/b^2+\cdots < \frac{b}{b-1}n
\]
buffer insertions, hence linear work for fixed $b$.  With
$k=\ceil{\log_b n}$ levels, at most $bk=O(b\log_b n)$ values are resident.
All buffers and promotions are bounded finite states, so the algorithm and its
resource counter are primitive recursive.  The remedian is an estimator, not
in general the exact sample median.

\section{Exact sampling distributions}

Let a finite law have ordered support $v_0<\cdots<v_{s-1}$ and weights
$w_j$, total $W$.  Draw $N=2n+1$ independent observations.  Independence is
represented by the finite product space, whose outcome weights are products
of the coordinate weights.

For a fixed $j$, put
\[
 A_j=\sum_{h<j}w_h,\qquad B_j=w_j,\qquad C_j=\sum_{h>j}w_h.
\]

\begin{theorem}[Exact odd-sample median mass]\label{thm:sample-pmf}
The numerator of $\Pr(M_N=v_j)$ over denominator $W^N$ is
\begin{equation}\label{eq:sample-pmf}
 \sum_{i=0}^{n}\sum_{k=0}^{n}
 \frac{N!}{i!\,k!\,(N-i-k)!}
 A_j^i C_j^k B_j^{N-i-k},
\end{equation}
where terms with $N-i-k\leq0$ are omitted.
\end{theorem}

\begin{proof}
The sample median equals $v_j$ exactly when at most $n$ observations are
strictly below $v_j$, at most $n$ are strictly above, and the remaining
$N-i-k$ observations equal $v_j$.  For fixed $i,k$, the multinomial
coefficient chooses the three classes, and the product of class weights gives
the displayed monomial.  The cases are disjoint and exhaustive.
\end{proof}

\begin{theorem}[Exact odd-sample CDF]\label{thm:sample-cdf}
Let $P_j=\sum_{h\leq j}w_h$.  Then
\begin{equation}\label{eq:sample-cdf}
 \Pr(M_N\leq v_j)
 =\frac1{W^N}\sum_{r=n+1}^{N}\binom Nr P_j^r(W-P_j)^{N-r}.
\end{equation}
\end{theorem}

\begin{proof}
The median is at most $v_j$ exactly when at least $n+1$ of the $N$
observations are at most $v_j$.  Their count has the finite binomial law with
success weight $P_j$.
\end{proof}

Both formulas use only factorials, powers, bounded sums, and integer
arithmetic.  They are direct \iPRA\ targets.

\subsection{Continuous density as an effective limit}

For a continuous CDF $F$, \cref{eq:sample-cdf} defines the polynomial
\[
 H_N(p)=\sum_{r=n+1}^{N}\binom Nr p^r(1-p)^{N-r}.
\]
Thus
\[
 \Pr(v<M_N\leq v+h)=H_N(F(v+h))-H_N(F(v)).
\]
If a regular density name supplies
\[
 F(v+h)-F(v)=f(v)h+R(v,h),
 \qquad \abs{R(v,h)}\leq\omega(\abs h)\abs h,
\]
with primitive-recursive $\omega(r)\to0$, a finite polynomial difference
identity yields the density
\[
 \frac{(2n+1)!}{n!n!}F(v)^n(1-F(v))^nf(v)
\]
with an explicit interval error.  This replaces informal differentials by a
finite difference and a modulus.

If $U=F(M_N)$ and $F$ is effectively continuous and strictly increasing, then
$U$ has the beta density
\[
 \frac{u^n(1-u)^n}{B(n+1,n+1)}.
\]
Finite beta-integral identities give
\[
 \mathbb E U=\frac12,
 \qquad \operatorname{Var}(U)=\frac1{4(2n+3)}
 =\frac1{4(N+2)}.
\]
These identities are polynomial/factorial calculations once integration of
polynomials on $[0,1]$ is represented by exact rational antiderivatives.

\subsection{Effective asymptotic normality}

The usual asymptotic theorem requires more than the mere existence of a
density.  Assume:

\begin{enumerate}
\item $m$ is a unique median;
\item $f(m)=c>0$ is a positive regular real with a rational lower bound;
\item a primitive-recursive local modulus certifies
$F(m+h)=1/2+ch+o(h)$;
\item the finite binomial CLT is available with an explicit error modulus.
\end{enumerate}

For fixed rational $z$, set $h_N=z/(2c\sqrt N)$ using regular square-root and
reciprocal names.  Then
\[
 F(m+h_N)=\frac12+\frac{z}{2\sqrt N}+o(N^{-1/2}).
\]
By \cref{thm:sample-cdf}, the event $M_N\leq m+h_N$ is a binomial tail.
Standardizing that binomial count and applying the effective CLT gives
\[
 2c\sqrt N\,(M_N-m)\Longrightarrow N(0,1),
\]
with an explicit rational error bound at rational thresholds.  Consequently
\[
 \operatorname{Var}(M_N)\sim\frac1{4Nc^2}
\]
when the required uniform-integrability modulus is supplied.

For a normal law of variance $\sigma^2$,
$c=1/(\sqrt{2\pi}\sigma)$, so the asymptotic median variance is
$\pi\sigma^2/(2N)$.  The mean variance is $\sigma^2/N$; their ratio tends to
$2/\pi$.  Every transcendental constant here is a regular name, and the limit
is an effective $\varepsilon$--$N$ theorem.

Delete-one jackknife inconsistency, delete-$k$ consistency, bootstrap
consistency, and the quoted $N^{-1/4}$ rate require a separate effective
resampling theory.  Their finite estimators are primitive recursive, but this
packet does not claim the asymptotic theorems.

\section{Derived estimators and regression}

\subsection{Hodges--Lehmann and pseudomedians}

For a sample $x_0,\ldots,x_{n-1}$, form the Walsh-average list
\[
 \left\{\frac{x_i+x_j}{2}:0\leq i\leq j<n\right\}.
\]
Its median is the one-sample Hodges--Lehmann estimator.  The population
pseudomedian is the median of $(X+X')/2$ for an independent copy $X'$.

If the original finite law is symmetric about $c$, reflection sends a pair
$(x_i,x_j)$ to a pair whose average is $2c-(x_i+x_j)/2$, preserving
multiplicity.  Hence the Walsh-average law is symmetric about $c$ and
\cref{thm:symmetry} makes $c$ a pseudomedian.  Positive affine equivariance is
immediate from the formula; negative affine maps use
\cref{thm:monotone}.

Robustness and asymptotic efficiency of the estimator are analytic extensions,
not consequences of this finite symmetry argument alone.

\subsection{Theil--Sen slopes}

For pairs with $x_i\neq x_j$, form
\[
 s_{ij}=\frac{y_j-y_i}{x_j-x_i}.
\]
The Theil--Sen slope is a median of this finite rational list.  It is invariant
under permutation of the observations.  Under
$x'=cx+d$, $y'=ay+b$ with $ac\neq0$, every slope becomes $(a/c)s_{ij}$, so
its median transforms accordingly.

If all points lie on $y=\alpha x+\beta$, every defined pairwise slope is
$\alpha$; therefore the Theil--Sen estimator recovers the exact slope.  This
is a one-line consequence of the median of a constant list and is fully
finitary.

\subsection{Median--median lines}

Sort bivariate observations by $x$, divide them into three consecutive groups,
and form the coordinatewise median point of each group.  If all observations
lie on an affine line $y=\alpha x+\beta$, affine monotone transport shows that
each group-median point lies on the same line, even when $\alpha<0$.  The line
through the outer group medians therefore recovers the original line; any
usual adjustment toward the middle median is zero.  General robustness claims
are algorithm-specific and are not asserted here.

\section{Multivariate medians}

\subsection{Marginal medians and the $\ell^1$ objective}

For points $x_i=(x_{i1},\ldots,x_{id})\in\Rat^d$, a marginal median is a
vector $m$ whose $j$th coordinate is a median of
$x_{1j},\ldots,x_{nj}$.

\begin{theorem}[Coordinatewise optimality]\label{thm:marginal}
The marginal median vectors are exactly the minimizers of
\[
  \sum_{i<n}\norm{x_i-a}_1.
\]
\end{theorem}

\begin{proof}
Expand the objective by coordinates:
\[
 \sum_i\norm{x_i-a}_1
 =\sum_{j<d}\sum_i\abs{x_{ij}-a_j}.
\]
The coordinates are independent optimization variables.  Apply
\cref{thm:l1} to each inner sum.
\end{proof}

This theorem is finite rational arithmetic and is directly internalizable for
fixed or coded dimension.

\subsection{Geometric median}

The geometric median minimizes
\[
 G(a)=\sum_i\norm{x_i-a}_2.
\]
Euclidean square roots require regular names, but several structural proofs
are constructive.

\begin{proposition}[Similarity equivariance]
If $T(a)=cRa+b$, where $c>0$ and $R$ preserves the Euclidean norm, then
$G_{T(x)}(T(a))=cG_x(a)$.  Thus $T$ carries geometric medians to geometric
medians.
\end{proposition}

\begin{proof}
For each $i$,
$\norm{T(x_i)-T(a)}_2=c\norm{R(x_i-a)}_2=c\norm{x_i-a}_2$.  Sum over $i$.
\end{proof}

\begin{theorem}[Correct uniqueness condition]\label{thm:geom-unique}
If the sample points are not all collinear, the geometric median is unique.
The hypothesis cannot be replaced merely by ``ambient dimension at least
two.''
\end{theorem}

\begin{proof}
Suppose distinct points $p,q$ both minimize $G$.  Convexity gives
\[
 G\left(\frac{p+q}{2}\right)
 \leq\frac{G(p)+G(q)}2=\min G,
\]
so equality holds.  Equality must hold term by term in the triangle
inequality
\[
 \norm{x_i-\frac{p+q}{2}}_2
 \leq\frac12\norm{x_i-p}_2+\frac12\norm{x_i-q}_2.
\]
Equality in the Euclidean triangle inequality forces $x_i,p,q$ to be
collinear.  Thus every sample point lies on the line through $p,q$, contrary
to the hypothesis.

For the counterexample to the unrestricted claim, take only
$(-1,0)$ and $(1,0)$.  Every point on the segment between them has total
distance $2$ and is a geometric median.
\end{proof}

Existence can be expressed effectively without compactness.  A minimizer lies
in the coordinate bounding box: clipping any coordinate to that box cannot
increase any distance.  For requested error $\varepsilon>0$, search a rational
grid of mesh $\delta\leq\varepsilon/(nd)$.  Moving by at most $\delta$ in each
coordinate changes each Euclidean distance by at most the $\ell^1$ movement,
so the total objective changes by at most $nd\delta\leq\varepsilon$.
Regular square-root intervals let a bounded search choose a certified
$\varepsilon$-minimizer.  This produces a primitive-recursive approximation
schema; an exact geometric median is the nested limit name.

The multivariate mean--spatial-median inequality follows by the same chain as
\cref{thm:mean-median}, with absolute value replaced by a norm, provided the
spatial median and norm are supplied by regular names.

\subsection{Centerpoints}

A centerpoint of $n$ points in $\Rat^d$ is a point such that every closed
halfspace containing it contains at least $n/(d+1)$ sample points.

\begin{theorem}[Centerpoint theorem, finite rational form]\label{thm:centerpoint}
For each fixed $d$, every nonempty finite rational point multiset has a
centerpoint.  A rational centerpoint can be chosen.
\end{theorem}

\begin{proof}
For every subset $S$ of more than $dn/(d+1)$ sample indices, consider its
convex hull $C_S$.  Any $d+1$ such subsets have a common sample index: each
complement has fewer than $n/(d+1)$ indices, so the union of the $d+1$
complements has fewer than $n$ indices.  Hence the corresponding convex hulls
intersect.

By finite Helly in $\Real^d$, all the convex hulls $C_S$ intersect.  Choose
$c$ in their intersection.  Suppose an open halfspace $H$ not containing $c$
contained more than $dn/(d+1)$ sample points.  Taking those indices as $S$
would give $c\in C_S\subseteq H$, a contradiction.  Therefore every open
halfspace not containing $c$ has at most $dn/(d+1)$ points, which is equivalent
to the centerpoint property for closed halfspaces through $c$.

All defining points are rational polytopes.  A nonempty finite rational
polytope has a rational point: enumerate a feasible basic subsystem and solve
it by determinant identities.  Thus $c$ can be rational.
\end{proof}

For fixed $d$, finite Helly can be formalized constructively through Radon's
lemma and determinant signs, or by a Farkas certificate.  Determinant bounds
give a finite search range for numerator and denominator.  This makes
\cref{thm:centerpoint} an \iPRA\ theorem schema for each fixed dimension, but
the required geometry library is not yet present in the supplied Lean
blueprint.

The ``median in all directions'' statement requires a separate projection and
subgradient development and remains Category C.

\section{Conditional medians}

Let a finite joint law on values $x$ and observations $y$ have natural weights
$w_{xy}$.  For each $y$ with positive fiber weight
$W_y=\sum_xw_{xy}$, choose a median $m_y$ of the conditional weighted list.

\begin{theorem}[Finite conditional median is Bayes $L^1$ optimal]\label{thm:conditional}
For every function $g$ on the finite $y$-support,
\[
 \sum_{x,y}w_{xy}\abs{x-m_y}
 \leq
 \sum_{x,y}w_{xy}\abs{x-g(y)}.
\]
\end{theorem}

\begin{proof}
For each fixed $y$, \cref{thm:weighted-l1} gives
\[
 \sum_xw_{xy}\abs{x-m_y}
 \leq\sum_xw_{xy}\abs{x-g(y)}.
\]
Sum these finitely many inequalities over $y$.
\end{proof}

Thus the article's conditional-median optimality is fully valid for finite
laws.  The characterization that a conditional median under additive Gaussian
noise is linear only for Gaussian priors is a deep analytic theorem and is not
proved here.

\section{Interpolated medians}

Suppose a median category has center $m$, width $w$, and contains $j>0$
observations.  Let $i$ observations lie below it and $k$ above it, so
$N=i+j+k$.  Assume the observations within the category are uniformly spread
across $[m-w/2,m+w/2]$.

\begin{theorem}[Interpolated-median identity]\label{thm:interpolated}
The interpolated median is
\[
 m_{\mathrm{int}}
 =m-\frac w2+w\frac{N/2-i}{j}
 =m+\frac w2\frac{k-i}{j}.
\]
It is equivalent to the CDF formula
\[
 m+w\left[\frac12-\frac{F(m)-1/2}{f(m)}\right].
\]
\end{theorem}

\begin{proof}
A fraction $(N/2-i)/j$ of the median category is needed to raise cumulative
mass from $i$ to $N/2$.  Starting at the lower category edge gives the first
formula.  Since $N=i+j+k$, elementary algebra gives
\[
 -\frac12+\frac{N/2-i}{j}=\frac{k-i}{2j}.
\]
For empirical proportions, $F(m)=(i+j)/N$ and $f(m)=j/N$; substitution gives
the same expression.
\end{proof}

For the article's $i=22$, $j=33$, $k=45$, $m=3$, $w=1$ example,
\[
 m_{\mathrm{int}}=3+\frac{45-22}{66}=\frac{221}{66}
 =3.34848\ldots,
\]
which rounds to $3.35$.  This is exact rational arithmetic in \iPRA.

\section{Median filters and clustering}

\subsection{Median filter}

\begin{theorem}[Majority interval theorem]\label{thm:majority-filter}
Let an odd window contain $n$ values.  If more than $n/2$ of them lie in an
interval $[a,b]$, then its median lies in $[a,b]$.  In particular, if more than
half equal $c$, the median equals $c$.
\end{theorem}

\begin{proof}
If the median were below $a$, then every one of the more-than-half inlier
values would be strictly above it, violating the upper balance count.  The
case above $b$ is symmetric.
\end{proof}

Therefore a median filter exactly removes salt-and-pepper corruption on a
constant odd window whenever fewer than half the pixels are corrupted.  More
general claims about image quality or edge preservation are empirical, not
unconditional theorems.

\subsection{One-dimensional and coordinatewise $k$-medians}

The correct $k$-medians objective is
\[
 \Phi(c_1,\ldots,c_k)=\sum_i\min_j\norm{x_i-c_j}_1.
\]
It minimizes within-cluster $\ell^1$ distance; it does not maximize distances
between cluster medians.

Given centers, assign each point to a nearest center using a deterministic tie
rule.  This cannot increase $\Phi$.  Given assignments, replace each center by
a coordinatewise median of its cluster.  By \cref{thm:marginal}, this also
cannot increase $\Phi$.  If centers are chosen from the finite set of input
coordinate values and ties are canonical, only finitely many states exist.
A strict lexicographic tie-break on equal-objective states prevents cycles, so
the alternating process terminates at a local optimum.  Every step and its
objective comparison is primitive recursive.

\section{Median-unbiased estimators}

Let a finite estimator $T$ have a weighted distribution depending on a
parameter $\theta$.

\begin{theorem}[Loss characterization]
The estimator is median-unbiased at $\theta$ exactly when $\theta$ minimizes
\[
 a\longmapsto \mathbb E_\theta\abs{T-a}.
\]
If it is mean-unbiased, so that $\mathbb E_\theta T=\theta$, then $\theta$
minimizes $\mathbb E_\theta(T-a)^2$.
\end{theorem}

\begin{proof}
The first assertion is \cref{thm:weighted-l1}.  For the second, expand
\[
 \mathbb E(T-a)^2
 =\mathbb E(T-\theta)^2+(a-\theta)^2
 +2(a-\theta)\mathbb E(\theta-T).
\]
The last term is zero under mean-unbiasedness, leaving a nonnegative square.
\end{proof}

Strictly monotone transformations preserve median-unbiasedness by
\cref{thm:monotone}.  Arbitrary one-to-one transformations do not.  For the
equally weighted values $0,1,2$, the median is $1$.  The injection
$0\mapsto0$, $1\mapsto100$, $2\mapsto1$ produces transformed values
$0,100,1$, whose median is $1$, not $100$.  Thus order preservation, not mere
injectivity, is the correct hypothesis.

\section{Formalization in Lean and ProveIt}

\subsection{Repository baseline}

The repository baseline inspected for this packet is commit
\texttt{df3ff8cc28be5eb713a14809a9df167e9aa94ab1}, with Lean toolchain
\texttt{v4.32.0}.  The most relevant checked modules are the core definitions,
order, bounded sums, pairing, search, permutations, division, and binomial
identities in \texttt{IPRADerivations}.  The PA list-coding development already
contains a division-free median relation, an ordinary-list specification,
existence, functionality, and examples.

\subsection{Two formalization layers}

The archive separates two layers.

\begin{enumerate}
\item \texttt{MedianFiniteBlueprint.lean} gives readable ambient definitions
for finite rational lists and laws, plus records of theorem obligations.  This
layer is useful for testing and for reusing Mathlib finite combinatorics.
\item \texttt{MedianObjectLanguageTarget.lean} states the actual endpoint.
Each public result must be a closed formula $\varphi$ with a term of type
\[
 ([]:\operatorname{List}(\operatorname{Fm}\ \operatorname{PTm}\ 0))
 \ \vdash[\mathrm{PRAAx}]\ \varphi.
\]
\end{enumerate}

An ambient theorem about \texttt{List} or \texttt{Real} is not a substitute
for the second layer.

\subsection{Dependency-ordered construction}

\begin{enumerate}
\item Define valid sequence codes, length, bounded projection, signed integer
and rational comparison.
\item Define less-than and greater-than counts by bounded sums.
\item Define compare-swap, insertion, and insertion sort.  Prove sortedness and
permutation certificates.
\item Define lower and upper middle projections and prove
\cref{thm:middle-interval}.
\item Derive odd/even formulas, existence, functionality of twice the median,
quantile names, and monotone transport.
\item Define absolute-deviation folds and formalize the pairwise lower-bound
proof.  Add the weighted expansion or direct prefix-mass proof.
\item Prove the cleared mean--median inequality through the finite
sum-of-squares identity.
\item Add rank perturbation, breakdown, and majority interval theorems.
\item Define finite laws and product samples; prove
\cref{thm:sample-pmf,thm:sample-cdf} using the binomial library.
\item Build the estimator applications by mapping a finite source list to a
finite derived list and invoking the core median theorem.
\item Add regular names and moduli only after the exact finite layer is
kernel-checked.
\end{enumerate}

The completed modules should be imported from
\texttt{IPRADerivations.lean}, audited in \texttt{Checks/Axioms.lean}, and
transported through the repository's \texttt{pra\_iff\_native} theorem.

\section{Exact executable validation}

The supplied Python checker is not a proof oracle.  It is a regression suite
using only integers and \texttt{fractions.Fraction}.  The recorded run passes
$1{,}662{,}266$ exact cases, summarized below.

\begin{longtable}{@{}p{0.31\textwidth}r p{0.47\textwidth}@{}}
\toprule
Check & Cases & Coverage\\
\midrule
\endhead
finite balance and $L^1$ & 791 & All multisets of size at most seven over
$\{0,1,2,3,4\}$\\
weighted medians and mean bound & 1,359 & Weights $0,\ldots,3$ on supports of
width at most five\\
rank perturbation & 1,118,480 & All pairs of samples of size at most five over
four values\\
exact odd sample law & 1,278 & Supports of size at most four, weights at most
three, $N=1,3,5$\\
ninther rank & 362,880 & Every permutation of nine distinct ranks\\
median-filter majority & 174,760 & All odd windows through size seven over four
gray levels\\
interpolated median & 660 & Exact algebra and the $221/66$ example\\
symmetry and skew & 1,793 & Symmetric samples and the positive-third-moment
counterexample\\
Theil--Sen collinear recovery & 35 & Seven rational points on each of thirty-five
lines\\
replacement breakdown & 230 & Sample sizes through twenty\\
\bottomrule
\end{longtable}

Floating point is used only to print a decimal approximation to $221/66$.
Every assertion is exact.

\section{What remains genuinely open in this packet}

The following article claims are not presented as completed \iPRA\ proofs:

\begin{itemize}
\item the precise continuous unimodal $\sqrt{3/5}$ and median--mode bounds;
\item delete-one jackknife inconsistency, delete-$k$ consistency, bootstrap
consistency, and their rates;
\item the Gaussian-noise linear conditional-median characterization;
\item a complete median-in-all-directions development;
\item a checked fixed-dimensional Helly/Radon library;
\item a fully compiled effective sample-median CLT and regular geometric
median library.
\end{itemize}

For each, the report identifies the missing effective hypotheses and a route
to a formal proof.  Omitting them from Category A is deliberate: qualitative
classical existence, an ambient real theorem, or a hidden appeal to choice
would not satisfy the requested \iPRA\ standard.

\section{Conclusion}

The median article is dominated by a very strong finitary core.  Once the
median is expressed by two bounded counts, the central middle-interval theorem
organizes nearly every elementary claim.  Pairing proves $L^1$ optimality;
weighted expansion proves the finite probability version; a sum-of-squares
identity proves the mean--median bound without division or square roots;
finite multinomial counting gives the exact sample law; and many statistical
applications reduce to constructing a derived finite list.

The analytic claims are not discarded.  They are reformulated as effective
theorem schemas whose inputs carry primitive-recursive regular names and
moduli.  This separation makes the trust boundary explicit and provides a
clear path from the present human proof to a complete Lean
\texttt{Prf PRAAx} development.

\appendix
\section{Object-language formulas}

The following schematic formulas are the principal targets after sequence
operations are represented by primitive-recursive terms.  Here $v$ is a valid
sequence code and $n=\operatorname{len}(v)$.

\paragraph{Balance median.}
\[
 \operatorname{Median}(v,t):\equiv
  2\operatorname{ltCount}(v,t)\leq n\ \wedge\
  2\operatorname{gtCount}(v,t)\leq n.
\]

\paragraph{Twice conventional median.}
\[
 \operatorname{TwiceMedian}(v,M):\equiv
 M=\operatorname{lowerMiddle}(\operatorname{sort}(v))
  +\operatorname{upperMiddle}(\operatorname{sort}(v)).
\]

\paragraph{Weighted finite law.}
For weight code $w$ and value code $a$,
\[
 2\sum_{j<s}w_j\ind[a_j<t]\leq\sum_{j<s}w_j,
 \qquad
 2\sum_{j<s}w_j\ind[t<a_j]\leq\sum_{j<s}w_j.
\]

\paragraph{Cleared mean--median inequality.}
\[
 (2S-MW)^2\leq4(WQ-S^2).
\]

\paragraph{Exact odd-sample mass.}
For $N=2n+1$,
\[
 W^N\Pr(M_N=v_j)=
 \sum_{i\leq n}\sum_{k\leq n}
 \operatorname{Mult}(N;i,N-i-k,k)
 A_j^iB_j^{N-i-k}C_j^k.
\]
All divisions in the displayed multinomial coefficient are implemented by the
already-proved exact factorial divisibility identity, or the coefficient is
computed by Pascal recursion.

\section{Primitive-recursive pseudocode}

\begin{lstlisting}[language={}]
ltCount(v,t):
  c := 0
  for i < len(v):
    c := c + if nth(v,i) < t then 1 else 0
  return c

insert(a, sorted):
  scan from the left;
  place a before the first entry >= a;
  return the new code and an index-permutation certificate

sort(v):
  out := empty
  for i < len(v):
    out := insert(nth(v,i), out)
  return out

medianInterval(v):
  s := sort(v)
  return (nth(s,(len(v)-1)/2), nth(s,len(v)/2))

oddMedianCDFNum(weights,j,n):
  N := 2*n+1
  A := sum_{h<=j} weights[h]
  W := sum weights
  return sum_{r=n+1}^N choose(N,r)*A^r*(W-A)^(N-r)
\end{lstlisting}

The loops are bounded by input-derived natural numbers.  Each can therefore be
compiled to primitive recursion and verified by quantifier-free induction.

\section{Source and provenance notes}

The source article was consulted on 16 November 2026.  Wikipedia is used as a
statement inventory, not as the authority for proofs.  The arguments in this
report are self-contained or explicitly labeled as effective schemas and
formalization routes.

The ProveIt repository was inspected through its current main branch at the
commit recorded above.  Relevant design references are the iPRA README and
derivation library, the PA list-coding median implementation, and the finite
CLT packet.  Exact paths and commit metadata are listed in
\texttt{provenance/SOURCES.md}.

\begin{thebibliography}{9}
\bibitem{wiki-median}
Wikipedia contributors, \emph{Median}, English Wikipedia,
\url{https://en.wikipedia.org/wiki/Median}, accessed 16 November 2026.

\bibitem{proveit-ipra}
V. Reshetnikov, \emph{ProveIt: Primitive Recursive Arithmetic},
\url{https://github.com/VladimirReshetnikov/ProveIt/tree/main/Logic/PrimitiveRecursiveArithmetic},
repository inspected 16 November 2026.

\bibitem{proveit-derivations}
V. Reshetnikov, \emph{IPRADerivations},
\url{https://github.com/VladimirReshetnikov/ProveIt/tree/main/Logic/PrimitiveRecursiveArithmetic/Derivations},
repository inspected 16 November 2026.

\bibitem{proveit-listcoding}
V. Reshetnikov, \emph{PA List Coding: Predicates and Aggregates},
\url{https://github.com/VladimirReshetnikov/ProveIt/tree/main/Logic/PeanoArithmetic/ListCoding},
repository inspected 16 November 2026.

\bibitem{proveit-clt}
V. Reshetnikov, \emph{Constructive central limit theorems in iPRA},
\url{https://github.com/VladimirReshetnikov/ProveIt/tree/main/Logic/PrimitiveRecursiveArithmetic/docs/proofs/probability/central_limit_ipra_2026-11-15},
repository inspected 16 November 2026.
\end{thebibliography}

\end{document}
